\title{Direct Preference Optimization: Your Language Model is Secretly a Reward Model}

\begin{abstract}
Although large-scale unsupervised language models (LMs) acquire extensive world knowledge and some capacity for reasoning, exerting precise influence over their behavior remains challenging owing to the fully unsupervised paradigm in which they are trained. Standard approaches for enhancing such controllability involve collecting human judgments on the comparative quality of model outputs and further fine-tuning the LM to favor these preferences, frequently utilizing reinforcement learning from human feedback (RLHF). However, RLHF typically entails a multi-stage and sometimes unstable pipeline, initially training a reward model that represents the annotated preferences, and subsequently optimizing the LM via reinforcement learning to increase the expected reward, while also constraining deviation from the pre-trained model. This work introduces a novel reward model parameterization for RLHF, permitting the analytic computation of the associated optimal policy. This insight enables the resolution of the canonical RLHF objective using a straightforward classification loss alone. The proposed approach, termed \textit{Direct Preference Optimization} (DPO), demonstrates stable behavior, high performance, and computational efficiency; notably, it obviates the need for language model sampling or extensive hyperparameter search during fine-tuning. Our empirical evaluation demonstrates that DPO can adapt LMs to human preference signals at least as effectively as existing approaches. Specifically, DPO fine-tuning surpasses PPO-based RLHF in managing the sentiment of generated outputs, and is at least comparable, if not superior, in tasks such as summarization and single-turn dialogue, all while offering a markedly simpler training procedure.
\end{abstract}

\section{Introduction}
Large-scale unsupervised language models (LMs) trained on vast corpora display a range of unexpected abilities~\citep{chowdhery2022palm, brown2020language, touvron2023llama,bubeck2023sparks}. Nevertheless, since the training data originates from humans with divergent intentions, values, and expertise, these models inevitably absorb a spectrum of behaviors and knowledge, not all of which are suitable for direct emulation. For instance, while an AI coding assistant should ideally \textit{recognize} typical programming errors to address them effectively, it is preferable for the model’s code generation to reflect the higher standard, albeit less frequent, of expert programming found within the training distribution. Likewise, though the model should \textit{be cognizant} of a widely held misconception (e.g., one believed by 50\% of individuals), it is undesirable for the model to perpetuate such errors in half of its responses regarding that topic. Therefore, distinguishing and eliciting the \emph{targeted outputs and conduct} from the LM’s extensive \textit{latent knowledge and skills} is essential for creating AI that is reliable, effective, and aligned with human intent~\citep{ouyang2022training}. While current approaches generally shape LMs through reinforcement learning (RL) to align with human judgment, we demonstrate that the RL-derived preference objective leveraged in these systems can be satisfied precisely through a straightforward binary cross-entropy loss, streamlining the process of learning from preferences.

Broadly, standard methods imbue language models with desired traits by fine-tuning on collections of human preference data that typify safe and useful model behaviors. This preference adaptation occurs after initial large-scale unsupervised pre-training on diverse textual data. The most intuitive means of preference learning involves supervised fine-tuning on examples of human-annotated ideal responses; however, the predominant methods rely on reinforcement learning from human (or artificial) feedback (RLHF/RLAIF; \citep{christiano2017deep,bai2022constitutional}). In RLHF frameworks, a reward model is constructed based on a set of annotated human preferences, and RL is employed to train the LM to prioritize responses rated highly by the reward model while preventing substantial divergence from the base model. Although RLHF approaches yield language models that excel in interactive and coding tasks, the RLHF workflow is significantly more intricate than pure supervised learning: it demands the training of several models, continuous sampling from the model policy during RL, and thus incurs notable computational overhead.

This work introduces a method for tuning language models to align with human preferences without the need for explicit reward modeling or reinforcement learning. We introduce \textit{Direct Preference Optimization (DPO)}, an approach that implicitly pursues the same reward maximization goal with a KL-divergence regularization as established RLHF frameworks, but offers greater simplicity both in implementation and training. Conceptually, DPO enhances the log probability assigned to favored responses relative to disfavored ones. Crucially, it adjusts this process with per-example importance weights, mitigating model collapse commonly observed when utilizing a straightforward probability ratio approach. As with prior algorithms, DPO grounds its framework in a theoretical preference model—such as the Bradley-Terry model (\cite{bradley1952rankanalysis})—which quantifies the agreement between a proposed reward function and observed human preference judgments. In contrast to existing RLHF pipelines, which rely on the preference model to guide reward model training followed by policy optimization using that model, DPO leverages a variable substitution to directly define the preference loss in terms of the policy. Using a collection of annotated human preferences regarding model completions, DPO employs a binary cross entropy loss, directly tuning the policy to approximate the optimal response distribution for an implicitly defined reward function tailored to the collected preferences.

The central innovation of this paper is the introduction of Direct Preference Optimization (DPO), a reinforcement learning-free strategy for preference-based language model training. Our empirical findings demonstrate that DPO matches or surpasses the performance of prevailing preference learning methods—including those utilizing PPO-based RLHF—on benchmark tasks like sentiment control, summarization, and dialogue, when evaluated on models up to 6B parameters.

As self-supervised language models have scaled, they have acquired the ability to tackle certain tasks in a zero-shot setting \citep{radford2019language} or by employing few-shot prompting \citep{gpt3,megatron,chowdhery2022palm}. Nonetheless, substantial improvements in downstream performance and alignment with user objectives are observed when these models undergo fine-tuning on datasets composed of instructions and human-generated completions \citep{mishra-etal-2022-cross,sanh2022multitask,chung2022scaling,thoppilan2022lamda}. This process, known as ‘instruction-tuning’, enhances the generalization capabilities of large language models (LLMs) to new instructions outside their original instruction-tuning datasets, thereby improving overall usability \citep{chung2022scaling}. Although instruction tuning has demonstrated considerable success, collecting \textit{relative} human evaluations of response quality is often more feasible than assembling large collections of expert-labeled examples. Building on this, later studies have further refined LLMs using datasets that capture human preferences, yielding improvements across domains such as translation \citep{kreutzer-etal-2018-reliability}, summarization \citep{stiennon2022learning,ziegler2020finetuning}, story generation \citep{ziegler2020finetuning}, and instruction compliance \citep{ouyang2022training,ramamurthy2023is}. The typical methodology in these studies involves first fitting a neural network-based reward function to align with the preference data, often operationalized through models like the Bradley-Terry model \citep{bradley1952rankanalysis}, and subsequently training the language model to optimize this reward via reinforcement learning techniques, notably REINFORCE \citep{williams1992reinforce}, proximal policy optimization (PPO; \cite{schulman2017proximal}), or their modifications \citep{ramamurthy2023is}. Another related stream of research uses LLMs already adapted for instruction following with human feedback to generate additional synthetic preference-labeled data, targeting specific aspects such as safety or non-harmfulness \citep{bai2022constitutional}, often guided only by weak supervision, such as a rubric written in natural language to direct the LLM’s labeling. These approaches represent a synthesis of research on reinforcement learning for language models \citep{Ranzato2015SequenceLT,paulus2018a,wu2018learning} and general frameworks for preference-based learning from humans \citep{christiano2017deep,kupcsik2018learning}. Despite the conceptual attractiveness of leveraging relative preferences from humans, the process of fine-tuning large language models through reinforcement learning presents significant practical hurdles; the present work introduces a theoretically grounded method for optimizing models based on relative preferences without reliance on reinforcement learning.

Research on policy learning from preferences outside linguistic applications has encompassed both bandit and reinforcement learning frameworks, resulting in the development of a range of methodologies. When preferences or rankings are used instead of explicit rewards in contextual bandit learning, the resulting framework is known as the contextual dueling bandit (CDB; \cite{yue2012karmed,dudik2015contextual}). Since CDBs operate without absolute reward signals, theoretical guarantees rely on identifying a \textit{von Neumann winner}—that is, a policy whose average performance outperforms or ties any competing policy at least 50\% of the time \citep{dudik2015contextual}. A key distinction, however, is that in the CDB framework, preference feedback is provided interactively during training, whereas in human preference learning scenarios, models are usually trained using a fixed offline dataset annotated with preference comparisons \citep{yan2022human}. In a similar vein, \textit{preference-based RL} (PbRL) methods operate by collecting binary preferences derived from an \textit{unknown} internal 'scoring' mechanism instead of reward signals \citep{BusaFekete2014,ruiz2023dueling}. PbRL encompasses various techniques, some of which leverage off-policy preference data, but these typically involve first constructing an explicit estimate of the underlying scoring or reward model, and then optimizing with respect to that model \citep{jain2013learning,BusaFekete2014,christiano2017deep,sadigh2017active,kupcsik2018learning}. By contrast, our approach performs policy optimization in a single stage, directly maximizing alignment with preference feedback.

\section{Preliminaries}\label{section:prelims}

We provide an overview of the RLHF pipeline described by \citeauthor{ziegler2020finetuning} (and subsequently by \citep{stiennon2022learning, bai2022training, ouyang2022training}), which typically comprises three primary steps: (1) supervised fine-tuning (SFT); (2) collection of preference feedback and training of a reward model; and (3) reinforcement learning for policy refinement.

\textbf{SFT}: The RLHF process initiates by adapting a pre-trained language model to the target task(s) (e.g., dialogue, summarization) using supervised learning with high-quality task-relevant datasets, yielding the model $\pi^\text{SFT}$. 

\textbf{Reward Modelling Phase}: During the reward modelling stage, the SFT-tuned model is provided with prompts $x$ to generate pairs of responses $(y_1, y_2)\sim \pi^\text{SFT}(y \mid x)$. These response pairs are then evaluated by human annotators, who indicate which of the two completions is preferred, recorded as $y_w\succ y_l \mid x$ (with $y_w$ and $y_l$ corresponding to the favored and less favored answers, respectively). The observed preferences are presumed to be the manifestation of an unobserved reward function $r^*(y, x)$. There exist multiple probabilistic frameworks for modeling such preference data, with the Bradley-Terry (BT) model \cite{bradley1952rankanalysis} commonly adopted (though the more general Plackett-Luce model \citep{plackett1975analysis, luce2012individual} can also be utilized when rankings of multiple responses are available). Under the BT model, the human-generated preference probability $p^*$ is expressed as:

\begin{equation}\label{eq:bradley-terry}
    p^*(y_1\succ y_2 \mid x)=\frac{\exp\left(r^*(x, y_1)\right)}{\exp\left(r^*(x, y_1)\right) + \exp\left(r^*(x, y_2)\right)}.
\end{equation}

Given a fixed dataset of comparative preferences $\mathcal{D}=\bigl\{x^{(i)}, y_w^{(i)}, y_l^{(i)}\bigr\}_{i=1}^N$ drawn according to $p^*$, one may introduce a parametrized reward function $r_{\phi}(x, y)$ and fit its parameters using maximum likelihood estimation. By casting this as a binary classification task, the loss to be minimized becomes the negative log-likelihood:

\begin{equation}\label{eq:reward_model}
    \mathcal{L}_R(r_{\phi}, \mathcal{D}) = -\mathbb{E}_{(x, y_w, y_l)\sim \mathcal{D}}\bigl[\log \sigma(r_{\phi}(x, y_w)- r_{\phi}(x, y_l))\bigr]
\end{equation}

where $\sigma$ denotes the logistic sigmoid function. For language models, $r_{\phi}(x, y)$ is typically initialized with weights from the SFT model $\pi^\text{SFT}(y \mid x)$, extending it by appending a linear head to the final transformer layer to yield a scalar reward prediction \cite{ziegler2020finetuning}. To promote stability in the reward estimation, normalization is applied such that the expected value of the reward over the dataset satisfies $\mathbb{E}_{x,y\sim \mathcal{D}}\left[r_\phi(x, y)\right] = 0$ for all $x$.

\textbf{RL Fine-Tuning Phase}: During reinforcement learning, the trained reward model serves as a feedback mechanism for guiding the language model. As established in earlier literature~\citep{jaques2017sequence, jaques2020human}, the optimization target is expressed as

\begin{equation}\label{eq:RL}
\max_{\pi_{\theta}}  \mathbb{E}_{x\sim \mathcal{D}, y\sim \pi_{\theta}(y \mid x)}\bigl[r_{\phi}(x, y)\bigr] - \beta\mathbb{D}_{\textrm{KL}}\bigl[\pi_{\theta}(y\mid x)\mid \mid \pi_\text{ref}(y\mid x)\bigr],
\end{equation}

where $\beta$ modulates the allowed divergence from a fixed reference policy $\pi_\text{ref}$, which is usually taken as the initial SFT model $\pi^\text{SFT}$. In practical settings, the policy $\pi_\theta$ is also initialized with the weights of $\pi^\text{SFT}$. The KL divergence constraint is crucial to prevent excessive divergence from regions where the reward model is well-calibrated, and it helps preserve diversity in model outputs, thereby avoiding collapse onto a single mode that receives high reward. Since the text generation process is inherently discrete, this objective is not differentiable end-to-end and therefore reinforcement learning algorithms are employed for optimization. The prevailing strategy \citep{ziegler2020finetuning, stiennon2022learning, bai2022training, ouyang2022training} reformulates the objective by defining the reward as ${r(x, y) = r_{\phi}(x, y) -\beta (\log \pi_{\theta}(y\mid x) - \log \pi_\text{ref}(y\mid x))}$ and optimizes it using Proximal Policy Optimization (PPO) \cite{schulman2017proximal}.

\section{Direct Preference Optimization}\label{sec:DPO}

Given the significant computational and methodological difficulties associated with applying reinforcement learning to large-scale language model fine-tuning, we aim to establish a streamlined alternative that utilizes preference data for direct policy learning. In contrast to prior RLHF strategies, which first estimate a reward model and then optimize with RL, our approach exploits a specialized parameterization for the reward function that allows derivation of its optimal policy analytically, removing the need for iterative RL procedures. As will be elaborated below, the central insight is to identify an explicit analytical relationship between the reward function and its corresponding optimal policy. This relationship enables us to re-express the objective from one over reward models to one directly over the policies themselves, effectively using a change-of-variables transformation.

\textbf{DPO Objective Derivation.} We begin with the general reinforcement learning objective presented in Eq.~\ref{eq:RL}, considering a broad reward function $r$. Building on existing literature~\citep{peters2007reinforcement, peng2019advantage, korbak2022reinforcement, go2023aligning}, one can readily derive that the optimizer of the reward maximization problem with a KL divergence constraint, as formulated in Eq.~\ref{eq:RL}, yields the following optimal policy structure:

\begin{equation}\label{eq:op_policy}
    \pi_r(y\mid x) = \frac{1}{Z(x)}\pi_\text{ref}(y\mid x)\exp\left(\frac{1}{\beta}r(x, y)\right),
\end{equation}

Here, the normalizing constant $Z(x) =\sum_{y}\pi_\text{ref}(y\mid x)\exp\left(\frac{1}{\beta}r(x, y)\right)$ serves as the partition function. A full derivation appears in Appendix \ref{app:derivation1}. Even if $r_{\phi}$ is used as a maximum likelihood estimator for the true reward $r^*$, computation of the partition function $Z(x)$ remains computationally intensive \citep{korbak2022reinforcement, go2023aligning}, presenting challenges for direct application of this formulation. Nevertheless, it is possible to manipulate Eq.~\ref{eq:op_policy} to represent the reward as a function of the optimal policy $\pi_r$, the reference policy $\pi_\text{ref}$, and the unknown partition function $Z(\cdot)$. By applying the logarithm to both sides of Eq.~\ref{eq:op_policy} and performing algebraic simplification, we obtain:

\begin{equation}\label{eq:main_eq}
    r(x,y) =\beta \log \frac{\pi_r(y\mid x)}{\pi_\text{ref}(y\mid x)} + \beta \log Z(x).
\end{equation}

This alternative parameterization can be applied to the ground-truth reward $r^*$ and its optimal policy $\pi^*$. Importantly, in the context of the Bradley-Terry model, preference depends only on reward differences between two outputs: ${p^*(y_1 \succ y_2 \mid x) = \sigma(r^*(x, y_1) - r^*(x, y_2))}$. Substituting the reparameterized expression for $r^*(x, y)$ from Eq.~\ref{eq:main_eq} into Eq.~\ref{eq:bradley-terry} eliminates the partition function, enabling human preference probabilities to be written solely in terms of the optimal policy $\pi^*$ and the reference policy $\pi_\text{ref}$. Consequently, under the Bradley-Terry framework, the optimal RLHF policy $\pi^*$ fulfills:

\begin{equation}\label{eq:objective}
    p^*(y_1\succ y_2 \mid x)=\frac{1}{1 + \exp\left(\beta \log \frac{\pi^*(y_2\mid x)}{\pi_\text{ref}(y_2\mid x)} - \beta \log \frac{\pi^*(y_1\mid x)}{\pi_\text{ref}(y_1\mid x)}\right)}
\end{equation}

Full details of this derivation are given in Appendix~\ref{app:derivation2}. While Eq.~\ref{eq:objective} employs the Bradley-Terry model, analogous expressions can be developed for more general settings, such as those described by Plackett-Luce models~\citep{plackett1975analysis, luce2012individual}, which are detailed in Appendix~\ref{app:plackett_luce_models}.

With the likelihood of human preferences now parameterized in terms of the optimal policy rather than rewards, we can establish a maximum likelihood estimation procedure for a policy parameterized by $\pi_\theta$. Drawing a parallel with the reward modeling framework (e.g., Eq.~\ref{eq:reward_model}), the associated objective for policy optimization is then defined as:

\begin{equation}\label{eq:optimum_model}
    \mathcal{L}_\text{DPO}(\pi_{\theta}; \pi_\text{ref}) = -\mathbb{E}_{(x, y_w, y_l)\sim \mathcal{D}}\left[\log \sigma \left(\beta \log \frac{\pi_{\theta}(y_w\mid x)}{\

In this approach, we learn an implicit reward by employing an alternative form of parameterization, wherein the optimal policy corresponds directly to $\pi_\theta$. Furthermore, as our methodology can be interpreted as fitting a reparametrized Bradley-Terry model, it inherits established theoretical guarantees, including consistency when the distribution of preference data meets appropriate conditions \cite{bong2022generalized}. We further explore and compare the theoretical underpinnings of DPO with related methods in Section~\ref{sec:theory}.

\textbf{Mechanistic perspective on the DPO update.} To better understand the DPO algorithm’s mechanism, it is insightful to examine the gradient of its loss $\mathcal{L}_\text{DPO}$. Specifically, the derivative with respect to parameters $\theta$ is given by:

\begin{multline*}\label{eq:gradient}
    \nabla_\theta \mathcal{L}_\text{DPO}(\pi_\theta;\pi_\text{ref}) = \\ -\beta\mathbb{E}_{(x, y_w, y_l) \sim \mathcal{D}} \bigg[\underbrace{\sigma(\hat{r}_\theta(x, y_l) - \hat{r}_\theta (x, y_w))}_\text{higher weight when reward estimate is wrong}\bigg[\underbrace{\nabla_\theta\log \pi(y_w \mid x)}_\text{increase likelihood of $y_w$} - \underbrace{\nabla_\theta\log\pi(y_l \mid x)}_\text{decrease likelihood of $y_l$}\bigg]\bigg],
\end{multline*}

with $\hat{r}_\theta(x, y) = \beta \log \frac{\pi_\theta(y \mid x)}{\pi_\text{ref}(y \mid x)}$ representing the implicit reward defined by the language model $\pi_\theta$ relative to the reference model $\pi_\text{ref}$ (see Section~\ref{sec:theory} for further details). Conceptually, this gradient encourages higher probability for preferred responses $y_w$ and suppresses the likelihood of less favored responses $y_l$. Critically, the update is weighted according to how much more the implicit reward $\hat{r}_\theta$ favors the dispreferred outputs, normalized by $\beta$—capturing both the extent of misordering by the model and the impact of the KL regularization. Empirical evidence highlights the significance of this weighting term: omitting it, as in an unweighted baseline, can lead to failure modes such as language model collapse (see Appendix Table~\ref{tab:unlikelihood_generations}).

\textbf{DPO procedure.} The typical DPO workflow consists of the following steps: 1) For each prompt $x$, generate completions $y_1, y_2 \sim \pi_\text{ref}(\cdot \mid x)$ and annotate them with human preferences, thereby forming the offline preference dataset $\mathcal{D} = \{x^{(i)}, y_w^{(i)}, y_l^{(i)}\}_{i=1}^N$; 2) Train the language model $\pi_\theta$ by minimizing $\mathcal{L}_\text{DPO}$, given the reference model $\pi_\text{ref}$, the dataset $\mathcal{D}$, and the chosen $\beta$ parameter. 
In real-world applications, it is preferable to make use of pre-existing publicly released preference datasets rather than generating new completions and collecting fresh human annotations. Because these datasets are produced using $\pi^\text{SFT}$, we set $\pi_\text{ref} = \pi^\text{SFT}$ when such a model is accessible. If $\pi^\text{SFT}$ is unavailable, $\pi_\text{ref}$ is initialized by maximizing the likelihood over preferred completions ${(x, y_w)}$, that is, setting ${\pi_\text{ref} = \argmax_{\pi}\mathbb{E}_{x, y_w \sim \mathcal{D}}\left[\log \pi(y_w \mid x)\right]}$. This approach reduces the impact of the distribution mismatch between the (unknown) true reference distribution and the practical $\pi_\text{ref}$ adopted by DPO. More information on implementation specifics and the chosen hyperparameters can be found in Appendix~\ref{app:implementation}.

\section{Theoretical Analysis of DPO}
This section provides deeper insights into the DPO approach, offers theoretical justification, and discusses the benefits of DPO in relation to the limitations seen in actor-critic methods used for RLHF, such as PPO~\cite{schulman2017proximal}.

\label{sec:theory}

\subsection{Your Language Model Is Secretly a Reward Model}
DPO eliminates the need for both explicit reward fitting and a reinforcement learning phase, achieving policy learning through a single maximum likelihood loss. Observe that the DPO loss in Eq. \ref{eq:main_eq} can be interpreted as a Bradley-Terry model, where the reward function is parameterized as $r^*(x, y) = \beta \log\frac{\pi^*_\theta(y \mid x)}{\pi_\text{ref}(y \mid x)}$. Optimizing our parameterized model $\pi_\theta$ under this framework corresponds, via a change of variables, to optimizing the reward model as described in Eq. \ref{eq:reward_model}. In what follows, we develop the underlying theory for this reparameterization, demonstrate that it does not restrict the expressiveness of reward models, and show it allows for the precise recovery of the optimal policy. We begin by formalizing an equivalence relation over reward functions.

\begin{definition}
We define two reward functions $r(x, y)$ and $r'(x, y)$ as equivalent if and only if there exists a function $f$ such that $r(x, y)-r'(x, y) = f(x)$.     
\end{definition}

It is straightforward to verify that this constitutes an equivalence relation, thereby dividing the set of reward functions into distinct equivalence classes. The following two lemmas can be established:

\begin{lemma}\label{lemma:same_prefrence} 
Within the Plackett-Luce, and specifically the Bradley-Terry, preference frameworks, any two reward functions belonging to the same equivalence class produce identical preference distributions.
\end{lemma}

\begin{lemma}\label{lemma:same_policy}
    Any pair of reward functions from a single equivalence class leads to the same optimal policy in the associated constrained RL setting.
\end{lemma}

The arguments supporting these statements are routine and are provided in Appendix \ref{app:lemma1}. The initial lemma refers to the commonly recognized issue of under-specification inherent in the Plackett-Luce model family \cite{plackett1975analysis}. As a consequence of this ambiguity, it becomes necessary to apply additional identifiability conditions to establish guarantees for maximum likelihood estimators derived from Eq. \ref{eq:reward_model} \cite{bong2022generalized}. The second lemma clarifies that any reward function in a given equivalence class will yield the same optimal policy; therefore, in terms of optimization, our focus is limited to identifying an arbitrary member from the corresponding optimal reward class. We provide a formal proof for the following theorem in Appendix~\ref{app:thm1}:

\begin{theorem}\label{thm:main}
    Given suitable regularity conditions, every reward class consistent with the Plackett-Luce (and specifically Bradley-Terry) model admits the parameterization ${r(x, y) = \beta \log \frac{\pi(y\mid x)}{\pi_\text{ref}(y\mid x)}}$ for some policy model $\pi(y\mid x)$ and fixed reference policy $\pi_\text{ref}(y \mid x)$.
\end{theorem}

\begin{sproof}
    Take any reward function $r(x, y)$ that defines an optimal policy $\pi_r(y \mid x)$, according to Eq. \ref{eq:op_policy}. We demonstrate that a reward function in the equivalence class of $r$ may be written using the aforementioned reparameterization. Define the mapping $f$ by  
\begin{equation}
    f(r; \pi_\text{ref}, \beta)(x, y) = r(x, y) - \beta\log\sum_{y}\pi_\text{ref}(y\mid x)\exp\left(\frac{1}{\beta}r(x, y)\right)
\end{equation}
This operator $f$ performs normalization by subtracting the log-partition function of $\pi_r$. Since the normalization term depends solely on $x$, $f(r; \pi_\text{ref}, \beta)(x, y)$ remains in the equivalence class of $r(x, y)$. Substituting $r$ with the expression from the right-hand side of Eq.~\ref{eq:main_eq}—which is valid for any $r$—yields $f(r; \pi_\text{ref}, \beta)(x, y) = \beta \log \frac{\pi_r(y\mid x)}{\pi_\text{ref}(y\mid x)}$. Thus, the mapping $f$ constructs a representative from the equivalence class of $r$ with the desired parameterization, confirming that the proposed reparameterization does not restrict the expressiveness of the reward model.
\end{sproof}

Theorem~\ref{thm:main} can also be interpreted as determining the specific reward function chosen by the DPO reparameterization within each equivalence class—namely, the reward function that fulfills:

\begin{equation}\label{eq:lag_p}
     \sum_{y}\underbrace{\pi_\text{ref}(y\mid x)\exp\left(\frac{1}{\beta}r(x, y)\right)}_{=\pi(y\mid x)\text{, using Thm.~\ref{thm:main} reparam.}} = 1,
\end{equation}

In other words, $\pi(y\mid x)$ defines a legitimate probability distribution, as it remains non-negative and its values sum to one. Further, in light of Eq.~\ref{eq:op_policy}, Eq.~\ref{eq:lag_p} corresponds to the partition function of the optimal policy derived from the reward function $r(x, y)$. A central idea underlying the DPO algorithm is the introduction of additional constraints on the otherwise undetermined Plackett-Luce (and, in particular, Bradley-Terry) class of preference models. These constraints preserve the diversity of admissible reward models while making the optimal policy in Eq.~\ref{eq:op_policy} analytically solvable for every prompt $x$.

\subsection{Instability of Actor-Critic Algorithms} Our framework also permits an analysis of potential instabilities that may arise in standard actor-critic approaches commonly adopted for RLHF, such as PPO. We consider the RLHF workflow, specifically focusing on the fine-tuning phase as described in Section \ref{section:prelims}. Connections emerge here with the control as inference perspective \cite{levine2018reinforcement} for the constrained reinforcement learning formulation presented in \ref{eq:RL}. Given a parameterized policy $\pi_{\theta}(y\mid x)$, the objective is to minimize $\mathbb{D}_{\text{KL}}[\pi_{\theta}(y|x) \mid \mid \pi^*(y\mid x)]$, where $\pi^*$ denotes the optimal policy arising from Eq.~\ref{eq:optimum_model} under the influence of reward $r_{\phi}(y, x)$. Through straightforward algebraic manipulations, this leads to the following objective function:

\begin{equation}\label{eq:AC}
    \max_{\pi_{\theta}}\mathbb{E}_{\pi_{\theta}(y\mid x)}\bigg[\underbrace{r_{\phi}(x, y) -\beta\log\sum_{y}\pi_\text{ref}(y\mid x)\exp\left(\frac{1}{\beta}r_{\phi}(x, y)\right)}_{f(r_{\phi}, \pi_\text{ref}, \beta)} - \underbrace{\beta\log\frac{\pi_{\theta}(y\mid x)}{\pi_\text{ref}(y\mid x)}}_{\text{KL}}\bigg]
\end{equation}

This objective mirrors the one targeted in earlier studies 
\citep{ziegler2020finetuning, stiennon2022learning, bai2022training, ouyang2022training}, utilizing a reward formulation for $r_{\phi}$ that aligns with the DPO framework. Within this context, the normalization factor in $f(r_{\phi}, \pi_\text{ref}, \beta)$ can be viewed as the soft value function corresponding to the reference policy $\pi_\text{ref}$. Although this component does not alter the optimal policy, omitting it may introduce significant variance into the policy gradient estimates, thereby destabilizing the learning process. One potential remedy involves parameterizing the normalization term through a learned value function, but this introduces further optimization challenges. Another strategy observed in preceding works involves reward normalization relative to a human completion baseline, which effectively serves as a single-sample Monte Carlo approximation for the normalization term. By contrast, the DPO reparameterization produces a reward function inherently free from dependence on baseline terms.

\section{Experiments}
We now empirically assess how effectively DPO can train policies grounded directly in preference data. Initially, we consider a controlled scenario in text generation, posing the question: how does DPO balance reward maximization and KL-divergence regularization with respect to the reference policy, particularly in comparison with prevalent preference optimization techniques such as PPO? Subsequently, we expand our analysis to encompass larger models and more complex RLHF applications, like summarization and conversational tasks. Remarkably, DPO consistently achieves performance on par with, or surpassing, established baselines—including RLHF implemented with PPO and the selection of the top-performing sample out of $N$ rollouts evaluated by a learned reward function—while requiring minimal hyperparameter tuning. Prior to discussing these empirical findings, we outline our experimental methodology; further specifics are provided in Appendix~\ref{app:exp_details}.

\textbf{Tasks.} We investigate three distinct open-ended text generation tasks in our experimental setup. Across all settings, algorithms learn a policy using a preference dataset, $\mathcal{D} = \bigl\{x^{(i)}, y_w^{(i)}, y_l^{(i)}\bigr\}_{i=1}^N$. In the \textbf{controlled sentiment generation} task, $x$ serves as a movie review prompt taken from the IMDb dataset \cite{maas-EtAl:2011:ACL-HLT2011}, and the policy is required to produce a continuation $y$ exhibiting positive sentiment. To allow systematic assessment, preference pairs for model generations are constructed using an existing sentiment classifier; specifically, a pair $(y_w, y_l)$ is favored such that $p(\text{positive} \mid x, y_w) > p(\text{positive} \mid x, y_l)$. For the SFT baseline, we further train GPT-2-large on the training split of IMDb reviews until the loss stabilizes (see App~\ref{app:sentiment_details} for more specifics). For the \textbf{summarization} task, $x$ corresponds to a Reddit forum post and the objective is to generate a summary $y$ covering the post’s core content. Consistent with earlier approaches, we employ the Reddit TL;DR summarization dataset \citep{volske-etal-2017-tl} supplemented with preference annotations collected by \citeauthor{stiennon2022learning}. Our SFT model is trained on human-generated Reddit summaries\footnote{\url{https://huggingface.co/CarperAI/openai_summarize_tldr_sft}} utilizing the TRLX \citep{leandro_von_werra_2023_7790115} RLHF pipeline. Notably, the human preference annotations by \citeauthor{stiennon2022learning} were drawn from model outputs of a related, albeit separately trained, SFT model. Lastly, in the \textbf{single-turn dialogue} scenario, $x$ consists of a user-initiated query, which may range from scientific inquiries to personal advice requests. Here, the policy’s task is to generate an informative and engaging reply $y$; for this, we utilize the Anthropic Helpful and Harmless dialogue dataset \citep{bai2022training}, containing 170,000 conversations between people and an automated assistant. Each exchange concludes with two candidate responses generated by a large, unspecified language model, with a human label indicating the preferred one. Unlike the previous tasks, there is no pre-existing SFT model for this setting; we consequently fine-tune an off-the-shelf language model exclusively on the preferred response completions to derive the SFT baseline.

\textbf{Evaluation.} We employ two distinct evaluation strategies in our experiments. For the controlled sentiment generation task, where the ground-truth reward function (i.e., a sentiment classifier) is accessible, we assess each algorithm by plotting its trade-off frontier between the obtained reward and KL-divergence from the reference policy. This enables a direct analysis of each method’s ability to optimize the reward under constraints. By contrast, in practical scenarios where the ground-truth reward function is unavailable, we assess algorithms based on their \textit{win rate} relative to a baseline, relying on GPT-4 as a surrogate for human judgment to evaluate both the quality of summaries and the helpfulness of dialogue responses in summarization and single-turn dialogue tasks, respectively. For summarization, the baseline corresponds to test set reference summaries, while in dialogue, it is the preferred response annotated in the dataset. Although prior work suggests that large language models may outperform traditional automated metrics as evaluators \citep{Chen2023ExploringTU}, we supplement our study with a human evaluation, as described in Sec.~\ref{sec:human-judgments}, to validate the use of GPT-4 in our assessments. The findings reveal that GPT-4’s decisions align closely with those of human raters, often matching or exceeding inter-annotator agreement among humans.

\textbf{Methods.} Beyond DPO, we systematically investigate several prominent methods for training language models to align with human preferences. For the summarization task, we utilize zero-shot prompting with \textbf{GPT-J} \citep{gpt-j}, and for dialogue, we employ 2-shot prompting using \textbf{Pythia-2.8B} \citep{biderman2023pythia}. We also compare to the \textbf{SFT} model and introduce \textbf{Preferred-FT}, which is produced by fine-tuning on supervised signals from the selected completion $y_w$—drawn either from the SFT model (in the sentiment and summarization settings) or from a generic language model (in dialogue tasks). Another pseudo-supervised approach considered is \textbf{Unlikelihood}~\citep{welleck2019neural}, which optimizes the policy to increase the probability of $y_w$ while actively decreasing the probability assigned to $y_l$, with an adjustable coefficient $\alpha\in[0,1]$ governing the influence of the ‘unlikelihood’ penalty. Additionally, we include \textbf{PPO} \citep{schulman2017proximal}, using rewards learned from preference data, alongside \textbf{PPO-GT}, an oracle method that leverages the ground-truth reward in the sentiment control experiments. For PPO-GT, we employ two variants: a standard implementation \cite{leandro_von_werra_2023_7790115} and a version with normalized rewards and tuned hyperparameters for enhanced results (these modifications are also applied to standard PPO runs with learned rewards). Lastly, the \textbf{Best of $N$} baseline samples $N$ completions from the SFT model (or from Preferred-FT in the dialogue setting) and selects the highest-reward response according to the learned reward function. While this baseline decouples reward model quality from PPO-based optimization and typically yields strong performance, its computational overhead is prohibitive for even moderate values of $N$ because it requires sampling $N$ responses per query during testing.

\subsection{How well can DPO optimize the RLHF objective?}

In standard RLHF methods, the KL-constrained reward maximization objective serves to balance reward exploitation with maintaining proximity to a reference policy. Thus, any algorithmic comparison must jointly consider both the attained reward and the extent of KL divergence; pursuing higher reward at the cost of a substantially larger KL divergence may not be advantageous. Figure~\ref{fig:frontier-tldr-main} illustrates the trade-off curve between reward and KL divergence for a variety of algorithms in the sentiment task. For each method, we conduct several independent training runs, each with a distinct setting for the policy conservativeness parameter (target KL $\in\{3,6,9,12\}$ for PPO, $\beta \in \{0.05,0.1,1,5\}$, $\alpha\in\{0.05,0.1,0.5,1\}$ for unlikelihood, and varying random seeds for preferred-FT), resulting in a total of 22 experiments. After every 100 training updates, policies are evaluated on a fixed set of test prompts, measuring both the mean reward with respect to the ground truth reward function and the average sequence-level KL\footnote{Specifically, this is computed as the total KL-divergence over all timesteps in the sequence.} relative to the reference policy $\text{KL}\left(\pi\mid \mid \pi_\text{ref}\right)$. The findings indicate that DPO yields the most favorable frontier, simultaneously achieving higher rewards and lower KL divergence compared to other approaches. This is particularly striking for several reasons. Firstly, both DPO and PPO target the same optimization objective, yet DPO attains superior efficiency; DPO provides a strictly better reward/KL profile compared to PPO. Secondly, even when PPO is provided access to the exact reward function (PPO-GT), DPO outperforms PPO along the frontier.

\subsection{Can DPO scale to real preference datasets?}
\label{sec:dpo-real-datasets}
We proceed to assess the fine-tuning capabilities of DPO on tasks such as summarization and single-turn dialogue. It is well established that automatic metrics like ROUGE often display weak alignment with human judgements in summarization~\citep{stiennon2022learning}, motivating earlier research that leverages human feedback through PPO to enhance summary quality. To compare various approaches, we generate outputs on the test split of the TL;DR summarization dataset and calculate the mean win rate relative to reference summaries from the test set. For each method, completions are drawn using sampling temperatures between 0.0 and 1.0; win rates are presented in Figure~\ref{fig:frontier-tldr-main} (right). The models—DPO, PPO, and Preferred-FT—all initiate from the identical GPT-J SFT checkpoint\footnote{\url{https://huggingface.co/CarperAI/openai_summarize_tldr_sft}} for fine-tuning. Our results indicate that DPO achieves an approximate win rate of 61\% at a temperature setting of 0.0, surpassing PPO, which reaches about 57\% at its optimal temperature (also 0.0). DPO additionally attains a greater peak win rate relative to the strongest performance among the $N$ baseline candidates. Importantly, DPO’s $\beta$ hyperparameter was not systematically optimized for these experiments, so the reported results may not reflect its full capabilities. Furthermore, DPO exhibits increased robustness across different sampling temperatures compared to PPO, which tends to lose performance—approaching that of the underlying GPT-J model—at elevated temperatures. Preferred-FT, by contrast, demonstrates minimal improvement over the SFT baseline. We further examine DPO versus PPO using human evaluations, as detailed in Section~\ref{sec:human-judgments}: at a temperature of 0.25, DPO-generated outputs were favored in 58\% of pairwise comparisons with PPO outputs at the same temperature.

For single-turn dialogue tasks, we conduct evaluations on a subset of the Anthropic HH dataset \citep{bai2022training} test split, focusing specifically on instances with a single human-assistant exchange. For GPT-4 based assessments, we use the test set's preferred completions as the benchmark references to determine each method's win rate. Since there is no canonical SFT baseline available for this setting, our process begins with a pre-trained Pythia-2.8B model. We apply Preferred-FT to train a reference model on selected completions, ensuring the completions remain in the model's distribution, and subsequently perform DPO training. Additionally, we benchmark against the best completion selected out of 128 Preferred-FT samples (our analysis indicates that the Best of $N$ baseline reaches saturation at 128 completions for this scenario; refer to Appendix Figure~\ref{fig:best-of-n}) and a 2-shot prompted Pythia-2.8B base model. Our findings reveal that DPO matches or surpasses performance for each method when evaluated with their optimal temperature settings. We further include an RLHF model trained using PPO on the Anthropic HH dataset\footnote{\url{https://huggingface.co/reciprocate/ppo_hh_pythia-6B}}, sourced from a widely recognized repository\footnote{\url{https://github.com/CarperAI/trlx/tree/main/examples/hh}}, but are unable to identify a prompt or temperature for this RLHF model that outperforms the original Pythia-2.8B. Drawing on results from TL;DR and acknowledging that both Best of 128 and PPO approaches are optimized for the same reward function, we use Best of 128 as an approximate stand-in for PPO-level outcomes. In summary, DPO emerges as the only computationally tractable method that reliably outperforms the preferred completions on the Anthropic HH dataset, delivering performance on par with or better than the much more resource-intensive Best of 128 approach. Lastly, Figure~\ref{fig:dialogue-main} illustrates the rapid convergence of DPO to its peak performance.

\subsection{Generalization to a new input distribution}

To further investigate how PPO and DPO handle distribution shifts, we test both sets of policies—trained during the Reddit TL;DR summarization experiment—on out-of-domain data comprising news articles from the CNN/DailyMail test split \citep{nallapati-etal-2016-abstractive}. For inference, we apply the optimal sampling temperatures identified in the TL;DR setting (0 and 0.25). The associated outcomes are displayed in Table~\ref{tab:ood}. We assess the win rate of each method by comparing its summaries to ground-truth references, leveraging GPT-4 with the same evaluation prompt used for TL;DR but substituting ``forum post'' with ``news article''. On this out-of-distribution evaluation, DPO maintains a substantial advantage over PPO. These findings indicate that DPO is capable of matching or exceeding PPO's ability to generalize, despite DPO not having access to the additional unlabeled Reddit TL;DR prompts utilized by PPO.

\subsection{Validating GPT-4 judgments with human judgments}
\label{sec:human-judgments}
To assess the trustworthiness of GPT-4-based evaluations, we conduct a human study on outputs from the TL;DR summarization experiment, employing two GPT-4 prompting strategies. The \textbf{GPT-4 (S)} prompt frames a straightforward request to select the summary that better conveys the essential information from the post. In contrast, the \textbf{GPT-4 (C)} prompt includes an explicit request to prefer more concise summaries; we investigate this because under the simple prompt, GPT-4 tends to favor verbose, repetitive responses more than human evaluators do. Complete prompts are available in Appendix~\ref{app:prompts}. We conduct three key comparisons to encompass a range of sample qualities: the best-performing approach (DPO, temp. 0.25), the lowest-performing (PPO, temp. 1.0), and a mid-level method (SFT, temp. 0.25), each against PPO sampled greedily at its optimal temperature. Our results reveal that, regardless of prompt, GPT-4's preferences correlate with human judgments to a similar extent as inter-human agreement, indicating GPT-4 serves as a reasonable surrogate for human evaluation. Due to limitations in the number of available raters, multiple human assessments are gathered only for DPO and PPO-1 comparisons. In sum, the \textbf{GPT-4 (C)} prompt more accurately mirrors human preferences; thus, we utilize this prompt for principal evaluations in Section~\ref{sec:dpo-real-datasets}. Further methodological specifics, such as the rater interface and participant roster, can be found in Appendix~\ref{app:human-study}.

\section{Discussion}
Learning from preferences offers a robust and scalable approach for developing capable and aligned language models. In this work, we presented DPO, a straightforward framework for preference-based training of language models that does not rely on reinforcement learning. Instead of reformulating preference learning to fit a traditional RL paradigm solely to utilize standard RL methodologies, DPO establishes a correspondence between language model policies and reward structures. This mapping allows for the direct optimization of a language model according to human preferences by applying a simple cross-entropy objective, obviating the need for reinforcement learning or a loss in general applicability. Notably, DPO matches or outperforms prevailing RLHF methods—including those leveraging PPO—without significant hyperparameter adjustment, thereby lowering the barrier for preference-driven language model training.

\textbf{Limitations \& Future Directions.} Several avenues for future exploration emerge from our findings. It remains to be thoroughly understood how DPO-trained policies fare in out-of-distribution generalization compared to those optimized with explicit reward signals; preliminary evidence indicates comparable generalization to models trained with PPO, but more exhaustive assessments are warranted. For instance, one question is whether self-labeling via DPO policies can harness unlabeled data as effectively. Moreover, it is an open question how phenomena such as reward over-optimization appear in the context of direct preference optimization, and whether the minor performance dip shown in Figure~\ref{fig:dialogue-main}-right is attributable to such effects. Our current experiments are limited to models up to 6B parameters, suggesting promising opportunities for extending DPO to significantly larger, state-of-the-art models. In terms of evaluation, we observe that the prompt used when querying GPT-4 can influence computed win rates; hence, optimizing prompt design for automated judgment could be a valuable subject for further inquiry. Beyond the scope of human preference learning for language models, DPO holds promise for training generative models across diverse modalities.